\documentclass[11pt]{article}
\usepackage[margin=1in]{geometry}
\usepackage{amsmath,amssymb}
\usepackage{enumitem}
\usepackage{xcolor}
\usepackage[colorlinks=true,linkcolor=blue,urlcolor=blue]{hyperref}
\setlength{\parskip}{0.4em}
\setlength{\parindent}{0pt}

\title{\vspace{-1.5em}Measuring the Dark-Matter Scale Height from Stellar Phase Space:\\
\large A Summary of the Simulation--Recovery Experiments}
\author{}
\date{\today}

\newcommand{\hdm}{h_{\rm dm}}
\newcommand{\Phider}{\Phi}

\begin{document}
\maketitle
\vspace{-3em}

\section{The aim of the project}

We want to measure the \textbf{dark-matter scale height} $\hdm$ of the Galactic
disk from the vertical phase-space distribution of stars, $(z, w)$ (height above
the midplane and vertical velocity), and---crucially---to understand \emph{when
and why} that measurement is biased.

The pipeline is a controlled, end-to-end numerical experiment:
\begin{enumerate}[nosep]
  \item \textbf{Generate} an initial stellar distribution function (DF) $f_0(z,w)$
        at time $t=0$.
  \item \textbf{Propagate} every star through a gravitational potential
        $\Phider(z)$ by Hamiltonian (leapfrog) flow to the observation time
        $t_{\rm obs}$, producing the ``observed'' phase-space cloud.
  \item \textbf{Recover} the potential parameters---chiefly $\hdm$---by
        \emph{back-integrating} the observed stars through a trial potential and
        maximising the likelihood of the implied initial conditions. The whole
        back-integration is differentiable, so the fit is done by gradient
        descent (PyTorch autograd).
\end{enumerate}

Because we generate the data ourselves, the ground truth is known exactly, so any
gap between the recovered and true $\hdm$ is a \emph{controlled bias} we can
dissect.

\paragraph{The key physical difficulty.} The dark-matter contribution to the
potential is
\[
  \Phi_{\rm DM}(z) = 4\pi G\,\rho_{\rm DM}\, \hdm^2\,\log\cosh\!\left(\frac{z}{\hdm}\right).
\]
Expanding near the midplane, the curvature $\Phi_{\rm DM}''(0)=4\pi G\rho_{\rm DM}$
is \textbf{independent of $\hdm$}; the scale height enters only through the
\emph{anharmonic} terms ($\sim z^4/\hdm^2$ and higher). Hence $\hdm$ is imprinted
only on high-$z$ / high-action stars and is intrinsically \emph{weakly
constrained}. This single fact explains most of the results below: anything that
perturbs the high-$z$ tail (a wrong DF shape, an unmodelled transient) can
masquerade as a change in $\hdm$, and a slowly time-varying $\hdm(t)$ is nearly
degenerate with a static end-point value.

\section{What each test does}

\paragraph{Test 1 --- Mismatched DF (a deliberate failure).}
Data are generated with one DF; the fit assumes a \emph{different} DF (e.g. a
mismatched velocity dispersion $\sigma$). This shows that an incorrect assumption
about the stellar distribution biases $\hdm$. A first attempt to ``patch'' the
mismatch with a small time-independent neural-network (NN) correction to the
potential is included---and shown to be insufficient.

\paragraph{Test 1b --- Bias asymmetry, quantified.}
A verification of the direction and size of the Test-1 bias using the leading
$r^2$ (harmonic) term and an exact-formula effective-radius fit. It confirms the
bias is a real, predictable consequence of the DF mismatch, not a numerical
artefact.

\paragraph{Test 2 --- Can a time-dependent potential remove the bias?}
One shared, biased dataset is fit with two different \emph{time-dependent}
corrections, swept over an amplitude cap $\alpha$: (a) a physically-motivated
\textbf{passing-satellite} model, and (b) a generic \textbf{Fourier $\hdm(t)$}.
\emph{Lesson:} extra time-dependent freedom removes the bias \emph{only if it has
the right functional form} (the satellite), not if it is a generic $\hdm(t)$. A
smooth $\hdm(t)$ barely changes the likelihood---the data are nearly blind to it
(the anharmonic/adiabatic degeneracy above)---so it cannot absorb a localised
transient.

\paragraph{Test 3 --- Give the fit an honest, larger freedom.}
The fit is handed the \emph{full} single-Gaussian DF ($\mu_z,\mu_w,\sigma_z,
\sigma_w$ all free) \emph{plus} a small NN potential bump capped at $10\%$ of the
static potential peak, on a \emph{clean} dataset (static halo, no satellite, truth
known exactly). Question: does this extra freedom still return the correct DF
($\mu\!\approx\!0$, $\sigma\!\approx\!$ truth), the correct $\hdm\!\approx\!250$~pc,
and a near-zero NN? This tests whether flexibility \emph{preserves} the answer
when there is no real feature to absorb.

\section{The two new tasks}

These extend Test 3 by relaxing the assumption that we know the shape of the
stellar DF---the honest stress test of the whole method.

\paragraph{Task 1 --- Double-Gaussian DF $+$ time-dependent potential
$\Phider(x)+\delta\Phi(x,t)$ \textnormal{\small(implemented)}.}
The initial DF becomes a \textbf{two-component (cold $+$ hot) Gaussian mixture}
with \emph{every} property free---both means and both dispersions of each
component, plus the mixing weight $\pi$---and the potential gains a
\textbf{time-dependent} correction $\delta\Phi(z,t)=A_{\rm nn}\,\mathrm{MLP}
(z/z_{\max},\,t/t_{\rm end})$, capped at $10\%$ of the static peak. On a clean
two-population dataset with a static halo, we ask whether the joint fit still
recovers the true two-population DF \emph{and} the correct $\hdm$, with
$\delta\Phi$ staying small. This checks that added DF freedom fixes DF-shape bias
without opening a new degeneracy with $\delta\Phi(x,t)$.

\paragraph{Task 2 --- Normalizing-flow DF $+$ $\Phider(x)+\delta\Phi(x,t)$
\textnormal{\small(planned)}.}
Task~1 pushed to the limit: $f_0$ becomes a \textbf{normalizing flow}, a neural
density model able to represent \emph{any} smooth initial distribution. The
dynamics are still forced to come entirely from the potential
$\Phider(x)+\delta\Phi(x,t)$---the flow describes only the initial condition. The
question ``\emph{does it still return the correct $\hdm$?}'' is genuinely
non-trivial: a maximally flexible DF could in principle make $\hdm$
unidentifiable. If $\hdm$ survives, the headline result follows: \textbf{the
dark-matter scale height can be recovered without assuming the stellar
distribution function}.

\section{The through-line}

Every test circles the same tension:
\begin{itemize}[nosep]
  \item Flexibility of the \emph{wrong shape} (a mismatched DF, a generic
        $\hdm(t)$) \textbf{biases} $\hdm$ or fails to fix an existing bias.
  \item Flexibility that is \emph{too generic} risks making $\hdm$
        \textbf{degenerate}.
  \item The goal is to find the freedom that is ``just right''---enough to absorb
        our ignorance of $f_0$, structured enough to leave the dynamical
        imprint of $\hdm$ intact.
\end{itemize}
Tasks~1 and~2 are the decisive experiments: they test whether letting $f_0$ be
flexible removes DF-shape bias \emph{without} creating a new degeneracy with the
time-dependent potential.

\end{document}
